\documentclass[11pt]{article}
\usepackage{amsmath,amsfonts}
\usepackage{graphicx}
%\usepackage{xcolor}
%\definecolor{gray}{rgb}{.75,.75,.75}
%\usepackage[landscape]{geometry}
%\usepackage{hyperref}
%\usepackage[bookmarks=true, pdfborder={0 0 .5}, urlbordercolor=gray]{hyperref}

\renewcommand{\thefootnote}{\fnsymbol{footnote}}


\addtolength{\topmargin}{-.8in}
\addtolength{\textheight}{1.4in}
\addtolength{\oddsidemargin}{-.4in}
\addtolength{\textwidth}{.8in}
\pagestyle{empty}

\newcommand{\ds}{\displaystyle}
\newcommand{\ts}{\textstyle}

\newcommand{\Z}{\mathbb{Z}}
\newcommand{\R}{\mathbb{R}}
\newcommand{\C}{\mathbb{C}}
\newcommand{\EE}{\mathcal{E}}
\newcommand{\tZ}{\tilde{\Z}}

\newcommand{\Sym}{\mathrm{Sym}}

\renewcommand{\circle}[1]{{\bigcirc\hspace{-.75em}#1}}

\begin{document}

\footnote[0]{Notes by Peter Magyar \ {\tt magyar@math.msu.edu}}
\vspace{-1em}

\centerline{\bf Math 133\hspace{1.5em} \hfill{\Large\bf
Arclength}\hfill Stewart \S8.1}
\vspace{1em}

\noindent 
{\bf Increments of length.} In this section, 
we give an integral formula to compute the length of a curve,
by the same Method of Slice Analysis we used in \S5.2 to compute volume,
and in \S5.3 to compute work (see end \S5.2).

We want the arclength $L$ of a graph curve $y=f(x)$ for $x\in[a,b]$. 
We cut the curve into $n$ bits determined by $\Delta x$-increments of $x\in[a,b]$.
(In the picture, $n=5$.)
\\[1em]
\centerline{
\includegraphics[width=3.6in]
{8-1-Length.png} 
}
\vspace{-.5em}

\noindent
Because the bit at the sample point $x_i$ is so short, 
it is well approximated by a straight segment, and we can 
use the Pythagorean Theorem to compute its length:
$$
\Delta L_i  \ \approx\ \sqrt{(\Delta x)^2+(\Delta y)^2}\,.
$$
We want to write this as a term in a Riemann sum, 
so we must write it in the form $g(x_i)\,\Delta x$ for some function $g(x)$.
We simply factor out $\Delta x$:
$$
\Delta L_i  
\ \approx\ 
\sqrt{\left(1+\tfrac{(\Delta y)^2}{(\Delta x)^2}\right)(\Delta x)^2}
\ =\ 
\sqrt{1+(\tfrac{\Delta y}{\Delta x})^2}\  \Delta x\,.
$$
In the limit as $n\to \infty$, we get $\Delta x\to0$  and $\frac{\Delta y}{\Delta x}\to \frac{dy}{dx} = f'(x_i)$, and the Riemann sum total of the $\Delta L_i$'s becomes an integral:
$$
L \ =\ \lim_{n\to\infty} \sum_{i=1}^n\Delta L_i \ =\ \lim_{n\to\infty}\sum_{i=1}^n\! \sqrt{1+(\tfrac{\Delta y}{\Delta x})^2}\ \Delta x
\ = \ \int_a^b\!\!\! \sqrt{1+(\tfrac{dy}{dx})^2}\ dx\,.
$$
In Newton notation:
$$
L  \ = \ \int_a^b\!\!\! \sqrt{1+y'(x)^2}\ dx\,.
$$
\\[.5em]
{\sc example:} Compute the arclength of the curve $y=x\sqrt x$ over the
interval $x\in[0,4]$. We have $\frac{dy}{dx} = (x^{3/2})' = \frac32x^{1/2}$,
so:
$$
L 
\ =\
\int_0^4 \sqrt{1+(\tfrac{dy}{dx})^2\,}\,dx
\ =\
\int_0^4 \sqrt{1+\tfrac94x\,}\,dx
\ =\
\left.\tfrac{8}{27} (1{+}\tfrac94x)^{3/2}\right|_{x=0}^{x=4}
\ =\
\tfrac{8}{27}(\sqrt{10}{-}1)
\ \approx\ 9.07
$$
To check this, we compare with the straight-line distance between the endpoints
$(0,0)$ and $(4, 8)$: this is $\sqrt{4^2+8^2} \approx 8.9$, and indeed the
length of the curve is slightly larger.
\\[.5em]
{\sc example:} Compute the circumference of the unit circle,
which is twice the arclength of the graph $y=\sqrt{1{-}x^2\,}$ for $x\in[-1,1]$:
$$
C  \ =\ 2L
\ = \ 
2\int_{-1}^1\!\! \sqrt{1+\left(\smash{\tfrac{d}{dx}\sqrt{1{-}x^2\,}}\right)^2}\ dx
\ = \ 
2\int_{-1}^1\! \sqrt{1+\left(\! \tfrac{-2x} {2\sqrt{1{-}x^2\,} }\!\right)^{\!\!2}}\ dx
$$
\\[-1em]
$$
\ = \ 
2\int_{-1}^1\! \sqrt{ \frac{1-x^2+x^2}{1-x^2} \,}\ dx
\ = \ 
2\int_{-1}^1\! \frac{1}{ \sqrt{1{-}x^2\,} }\ dx
\ =\
\left.\vphantom{\sum}2\arcsin(x)\,\right|_{x=-1}^{x=1}
\ =\
2\pi\,.
$$
\\[.5em]
{\sc example:} Compute the arclength of the parabola $y=x^2$ over any
interval $x\in[0,b]$.
$$
L 
\ =\
\int_0^b\!\!\! {\ts \sqrt{1+(\tfrac{d}{dx}(x^2))^2}}\ dx
\ =\ 
\int_0^b\!\!\! {\ts \sqrt{1+4x^2\,}}\ dx\,.
$$
To find the indefinite integral, we use the reverse trig substitution (\S7.3):
$x = \frac12\tan(\theta)$, \ $\sqrt{1+4x^2\,} = \sec(\theta)$, \ $dx = \frac12\sec^2(\theta)\,dx$:
$$
\int\!\! {\ts \sqrt{1+4x^2\,}}\ dx
%\ =\
%\int_0^b\! \sec(\theta)\,\cdot\,\sec^2(\theta)\,d\theta
\ =\
\int\! \tfrac12\sec^3(\theta)\,d\theta
\ =\
\tfrac14
\ln\!\left|\vphantom{A^2}\!\tan(\theta){+}\sec(\theta)\right| 
\,+\,\tfrac14\tan(\theta)\sec(\theta),
$$
where we use $\int\!\sec^3(\theta)\,d\theta$ from \S7.2.
Restoring the original variable, $\tan(\theta)=2x$, $\sec(\theta)=\sqrt{1{+}4x^2\,}$,
and taking the definite integral:
$$
L
\ =\
\left[
\tfrac14\ln\!\left|\vphantom{A^2} 2x{+}\sqrt{1{+}4x^2\,}\,\right| 
\,+\,\tfrac12x\sqrt{1{+}4x^2\,}\,
\right]_{x=0}^{x=b}
\ =\
\tfrac14\ln\!\left|\vphantom{A^2} 2b{+}\sqrt{1{+}4b^2\,}\,\right| 
\,+\,\tfrac12b\sqrt{1{+}4b^2\,}.
$$
Arclength tends to get quite complicated even for quite simple curves!
\\[1em]
{\sc example:} Compute the arclength of the curve $y=x^3$ over  $x\in[0,1]$.
$$
L 
\ =\
\int_0^1\!\!\! {\ts \sqrt{1+(\tfrac{d}{dx}(x^3))^2}}\ dx
\ =\ 
\int_0^1\!\!\! {\ts \sqrt{1+9x^4\,}}\ dx\,.
$$
This is already complicated enough that it has no algebraic antiderivative.\footnote{
The integral can be expressed in terms of an ``elliptic function'', but this is circular reasoning
since elliptic functions themselves are defined as integrals!}

Does this mean the arclength formula is useless? Not at all! We cannot get an answer
on the algebraic level, but we can still get a numerical answer as accurate as we like.
This means going from the integral formula for $L$ back to the Riemann sums from which we deduced the integral.  For example, taking $n=1000$, the increment is $\Delta x=\frac1{1000}=0.001$, and the sample points are $x_i=i\,\Delta x = (0.001)i$.
The computer gives:
$$
L
\ \approx\ 
\sum_{i=1}^{n} \sqrt{1+9x_i^4\,}\, \Delta x
\ =\
\sum_{i=1}^{1000} \sqrt{1+9(10^{-12})i^4\,}\, (0.001)
\ \approx\ 1.548\,,
$$
To gauge the accuracy of this, we re-do it with $n=10,000$, getting $L \approx 1.547$,
so we can be confident that $L \approx 1.54$ is accurate to 2 decimal places.



\end{document}

%%%%%%%%%%%  Picture  %%%%%%%%%%
\\[0em]
\centerline{
%\includegraphics[width=2in]
{1-8-Discont-iv.png}
\qquad \qquad 
%\includegraphics[width=2in]
{1-8-Discont-v.png} 
}
\vspace{0em}

\noindent
%%%%%%%%%%%%%%%%%%%%%%%%%

%%%%%%%%%%%  Table  %%%%%%%%%%
In fact, we get the following derivatives:
$$\begin{array}{|c||c|c|c|c|c|c|}
\hline &&&&&& \\[-.9em] 
f(x) & \sin(x) & \cos(x) & \tan(x)  & \sec(x) & \csc(x) & \cot(x)
\\[.2em] \hline &&&&&& \\[-.9em]
f'(x) &\cos(x) & -\sin(x) & \sec^2(x) & \tan(x)\sec(x) & -\cot(x)\csc(x) & -\csc^2(x)
\\[.1em] \hline
\end{array}$$
\\[1em]
%%%%%%%%%%%%%%%%%%%%%%%%%%



















